% ============================================================
% ch17a.tex —— 第 17 章前半（17.1–17.2.2）
% 底本：Springer LNM 2386, printed pp.205–226 = PDF pp.245–266
% 公式唯一真值：work/pages/pNNN.png（200 dpi）
% 本块只在章首写一次 \chapter；后半（17.2.3 起）由 ch17b 续写
% ============================================================

% [原文 p.205]
\chapter{四维超对称量子场论}

\ChapterAbstract{本章聚焦于带 $\mathcal{N}=1$ 超对称的 $4d$ 规范理论的一般框架，$\mathcal{N}=2$ 理论构成其特殊子集。我们首先讨论支配 $4d$ 量子场论的对称性，为理解超对称模型提供必要的背景。随后我们探讨 $4d$ $\mathcal{N}=1$ 规范理论的经典性质，特别强调 $\mathcal{N}=1$ 拉格朗日密度的构造。最后我们处理这些理论的量子方面，深入讨论重整化、反常、全纯性以及模空间的量子修正等关键课题。这些准备将使我们对 $\mathcal{N}=2$ 动力学乃至谱网络作更深入的研究。}

对于希望更深入钻研四维超对称的读者，有丰富的教材 [358, 136, 103, 355, 349, 341] 与讲义 [253, 262, 56, 22, 335, 21, 111, 50] 可提供易懂的介绍、详细的讨论以及大量应用。

\section{对称性}

\subsection{超庞加莱代数}

闵可夫斯基时空 $\mathcal{M}_4=(\R^{1,3},\eta_{\mu\nu})$ 中任何四维（$4d$）量子场论\footnote{我们只考虑局域且幺正的量子场论。}（QFT）的对称性都包含庞加莱群，后者是时空平移与洛伦兹群 $\mathrm{O}(1,3)$ 的半直积。更确切地说，Coleman–Mandula 定理断言：若 $G$ 是某个满足若干假设（例如粒子数有限与弱弹性解析性）的非平凡幺正相对论性 $S$-矩阵的对称性李群，则 $G$ 是庞加莱群与某个内部对称性群的直积 [94, 355]。

% [原文 p.206]
洛伦兹群有四个连通分支。含恒等元的分支称为限制洛伦兹群，记作 $\SO(1,3)^+$。商群 $\mathrm{O}(1,3)/\SO(1,3)^+$ 同构于 $\Z_2\times\Z_2$，对应于离散对称性 $1,\mathcal{P},\mathcal{T},\mathcal{PT}$。这里 $\mathcal{P}$ 表示宇称（空间反演），$\mathcal{T}$ 表示时间反演，而 $\mathcal{PT}$ 是宇称与时间反演的组合。在无穷小对称性的层面，有李代数同构
\[
\so(1,3)\simeq\su(2)\times\su(2)^*\,,
\]
其中上标 $*$ 表示复共轭交换两份 $\su(2)$。限制洛伦兹群的万有覆叠是自旋群 $\operatorname{Spin}(1,3)^+$，作为实李群它同构于 $\SL_2(\C)\simeq\Sp_2(\C)$。它是 $\SO(1,3)^+$ 的二重覆叠，即 $\SO(1,3)^+$ 同构于 Möbius 群 $\PSL_2(\C)$。群 $\SL_2(\C)$ 有两个不等价（实）的复二维不可约表示，对应于 $\so(1,3)$ 的两个不等价旋子\footnote{关于旋子及其与洛伦兹群关系的入门介绍，可参见例如 [333]。}表示。第一个表示中元素的各分量记为 $\psi_\alpha\in\C$，$\alpha=1,2$；第二个表示中记为 $\psi_{\dot\alpha}\in\C$，$\dot\alpha=\dot 1,\dot 2$（这就是 van der Waerden 记法）。矩阵 $A\in\SL_2(\C)$ 作用为
\[
\begin{pmatrix}\psi_1\\ \psi_2\end{pmatrix}\to A\begin{pmatrix}\psi_1\\ \psi_2\end{pmatrix},\qquad
\begin{pmatrix}\psi_{\dot 1}\\ \psi_{\dot 2}\end{pmatrix}\to A^*\begin{pmatrix}\psi_{\dot 1}\\ \psi_{\dot 2}\end{pmatrix}.
\]
$\psi_\alpha$（相应地 $\psi_{\dot\alpha}$）称为左手（相应地右手）外尔旋子，在厄米共轭下满足 $\psi_\alpha^\dagger=\psi_{\dot\alpha}$。

相应地，$\SL_2(\C)$ 以两种不同方式作用在 $\mathcal{M}_4$ 上，称为外尔表示。设 $(\sigma_\mu)=(\sigma_0,\sigma_1,\sigma_2,\sigma_3)$，其中 $\sigma_0=\mathrm{Id}_{2\times 2}$，诸 $\sigma_i$ 为通常的泡利矩阵：
\[
\sigma_1=\begin{pmatrix}0&1\\ 1&0\end{pmatrix},\qquad
\sigma_2=\begin{pmatrix}0&-\ii\\ \ii&0\end{pmatrix},\qquad
\sigma_3=\begin{pmatrix}1&0\\ 0&-1\end{pmatrix},
\]
并令 $(\sigma^\mu)=(\sigma_0,-\sigma_1,-\sigma_2,-\sigma_3)$。点 $x=(x^\mu)\in\R^{1,3}$ 可用两种方式改写为厄米特 $2\times 2$ 矩阵：
\[
X=x^\mu\sigma_\mu\,,\qquad \bar X=x_\mu\sigma^\mu\,,
\]
矩阵 $A\in\SL_2(\C)$ 在其上通过共轭作用：
\[
X\to AXA^\dagger\,,\qquad \bar X\to (A^{-1})^\dagger\bar XA^{-1}\,.
\]
这些映射描述了洛伦兹群的两个独立不可约手征旋子表示：$X$（相应地 $\bar X$）对应于左手（相应地右手）旋子。矩阵 $\sigma_\mu$（以及 $\sigma^\mu$）自然带有一个无点与一个有点指标 $\sigma^\mu_{\alpha\dot\alpha}$，由此可以构造不变量，例如：

% [原文 p.207]
\[
\psi\sigma^\mu\bar\chi:=\psi^\alpha\sigma^\mu_{\alpha\dot\alpha}\bar\chi^{\dot\alpha}\,.
\]
由于 $\SL_2(\C)\simeq\Sp_2(\C)$，存在一个反对称不变张量：
\[
\epsilon_{\alpha\beta}=\epsilon_{\dot\alpha\dot\beta}=\begin{pmatrix}0&-1\\ 1&0\end{pmatrix}
\qquad\text{与}\qquad
\epsilon^{\alpha\beta}=\epsilon^{\dot\alpha\dot\beta}=\begin{pmatrix}0&1\\ -1&0\end{pmatrix},
\]
借助它，例如有 $\psi^\alpha=\epsilon^{\alpha\beta}\psi_\beta$，并且它诱导一个标量积：$\psi\chi=\psi^\alpha\chi_\alpha=\chi\psi$。

在外尔表示中，狄拉克矩阵写为
\[
\gamma^\mu=\begin{pmatrix}0&\sigma^\mu\\ \bar\sigma^\mu&0\end{pmatrix},\qquad
\gamma_5=\begin{pmatrix}1&0\\ 0&-1\end{pmatrix},
\]
而一个狄拉克旋子分解为 $(\psi_\alpha,\bar\chi^{\dot\alpha})$。一个狄拉克旋子是马约拉纳旋子（即实的）当且仅当 $\psi=\chi$。

在某些 QFT 中，存在超出 Coleman–Mandula 定理所允许范围的额外对称性。出现这种情况的一个重要情形是：理论的对称性构成超李群而非通常的李群，即在玻色型生成元之上还有费米型（奇）对称性生成元。相应的无穷小对称性代数成为 $\Z_2$-分次的：
\[
\mathfrak{g}=\mathfrak{g}_0\oplus\mathfrak{g}_1\,,
\]
其中 $\mathfrak{g}_0$ 是偶子代数，含玻色型生成元；$\mathfrak{g}_1$ 是奇部，含费米型生成元。与物理约束（例如包含超对称）相容的超李代数由 Haag–Lopuszanski–Sohnius 定理 [182] 支配，后者把 Coleman–Mandula 的结果推广到包含这些费米型对称性的情形。

\begin{definition}[17.1] 设 $\mathcal{N}\in\N_{>0}$。$4d$ 中 $\mathcal{N}$-扩张超庞加莱代数由通常的庞加莱代数玻色型生成元以及奇生成元 $Q^I_\alpha$ 与 $\bar Q^I_{\dot\alpha}$（$I=1,\dots,\mathcal{N}$）生成，它们满足：

% [Tutor Note: 定义 17.1 的指标约定可对照 p.206–207：$\alpha,\beta$ 为无点（左手）指标，$\dot\alpha,\dot\beta$ 为有点（右手）指标；$\sigma^\mu_{\alpha\dot\alpha}$ 恰有一个无点与一个有点指标，$\bar\sigma^\mu$ 则为其转置。核对 (17.1) 的量纲/结构时注意：反对易子 $\{Q,\bar Q\}$ 给出动量 $P_\mu$，而 $\{Q,Q\}$ 给出中心荷 $\mathcal{Z}^{IJ}$，两类量纲相同（质量）。]
% [原文 p.208]
\begin{equation}
\begin{aligned}
[P_\mu,Q^I_\alpha]&=0\,,\\
[P_\mu,\bar Q^I_{\dot\alpha}]&=0\,,\\
[M_{\mu\nu},Q^I_\alpha]&=\ii(\sigma_{\mu\nu})_\alpha^{\ \beta}Q^I_\beta\,,\\
[M_{\mu\nu},\bar Q^{I\dot\alpha}]&=\ii(\bar\sigma_{\mu\nu})^{\dot\alpha}_{\ \dot\beta}\bar Q^{I\dot\beta}\,,\\
\{Q^I_\alpha,\bar Q^J_{\dot\beta}\}&=2\sigma^\mu_{\alpha\dot\beta}P_\mu\delta^{IJ}\,,\\
\{Q^I_\alpha,Q^J_\beta\}&=2\sqrt{2}\epsilon_{\alpha\beta}\mathcal{Z}^{IJ}\,,\\
\{\bar Q^I_{\dot\alpha},\bar Q^J_{\dot\beta}\}&=2\sqrt{2}\epsilon_{\dot\alpha\dot\beta}\left(\mathcal{Z}^{IJ}\right)^*\,.
\end{aligned}
\tag{17.1}
\end{equation}
\end{definition}

这里 $P_\mu$ 与 $M_{\mu\nu}$ 是时空平移与无穷小洛伦兹变换的通常玻色型生成元，并且
\[
(\sigma^{\mu\nu})_\alpha^{\ \beta}=\left(\sigma^\mu_{\alpha\dot\gamma}(\bar\sigma^\nu)^{\dot\gamma\beta}-\sigma^\nu_{\alpha\dot\gamma}(\bar\sigma^\mu)^{\dot\gamma\beta}\right)/4\,.
\]
整数 $\mathcal{N}$ 是代数的超对称个数；超代数中超荷 $Q^I_\alpha$ 与 $\bar Q^I_{\dot\alpha}$ 的数目为 $4\mathcal{N}$。生成元 $\mathcal{Z}^{IJ}$ 是中心元，称为中心荷；$\mathcal{Z}^{IJ}$ 对其指标反对称。此外，我们要求存在一个 $\SU(\mathcal{N})_R$ 对称性，按指标 $I$ 转动超荷（该记号将在 17.1.3 节解释），即所有 $\mathcal{N}=1$ 子代数都是等价的。

带有 $\mathcal{N}$ 个超对称的超庞加莱代数指数化为超庞加莱群 $\mathrm{OSp}(4|\mathcal{N})$。$\mathcal{N}=1$ 超空间是齐性空间
\[
\mathcal{M}_{4|1}=\mathrm{OSp}(4|1)/\SO(1,3)\,,
\]
其通常仿射坐标为 $(x^\mu,\theta_\alpha,\bar\theta_{\dot\alpha})$，$\mu=0,1,2,3$，$\alpha,\dot\alpha=1,2$，其中 $\theta$ 与 $\bar\theta$ 是反交换的 Grassmann 数。

这一讨论可推广到任意维数，得到带 $\mathcal{N}$ 个超对称的 $d$ 维庞加莱超代数。旋子表示本质上依赖于维数；特别地，其性质依赖于维数模 8（Bott 周期性）。

我们将主要关注仅含自旋/螺旋度小于 1 的基本粒子的 QFT（即非引力量子场论）。这一要求给出超对称数目 $\mathcal{N}$ 的上界，且该上界依赖于 $d$。例如当 $d=4$ 时，上界为 $\mathcal{N}\le 4$。

\subsubsection*{超对称代数的表示}

量子力学的一条公设是：量子系统联系着一个态希尔伯特空间 $\mathcal{H}$，系统的（纯）态与 $\mathcal{H}$ 中的复直线一一对应，即与射影空间

% [原文 p.209]
$\mathbb{P}(\mathcal{H})$ 的元素一一对应。$\mathcal{H}$ 中的复直线与秩一投影算符（纯态的密度矩阵）一一对应。两态之间的跃迁振幅由 $\mathcal{H}$ 中的内积给出，而系统的对称性须保持跃迁振幅。维格纳定理指出：量子系统的任何对称性——定义为射线程空间 $\mathbb{P}(\mathcal{H})$ 的等距——都可用 $\mathcal{H}$ 上的线性幺正或反线性反幺正算符表示。

\begin{definition}[17.2] 在 QFT 中，单粒子态或单粒子多重态是时空对称性群的幺正有限维不可约表示。\end{definition}

对一般的 $4d$ QFT，时空对称性群是四维庞加莱群。相应的代数有两个卡西米尔算符：$P^2$ 与 $W^2$，其中 $P_\mu$ 生成时空平移，而 $W^\mu=\frac{1}{2}\epsilon^{\mu\nu\rho\sigma}P_\mu M_{\rho\sigma}$ 是 Pauli–Lubanski 矢量。有质量粒子由其质量（$P^2=m^2$）与自旋 $j$ 确定。此时有 $W^2=-m^2j(j+1)$。无质量粒子满足 $P^2=W^2=0$，并在 $P^\mu=(E,0,0,E)$ 的参考系中有 $W^\mu=\pm jP^\mu$，其中 $j$ 是螺旋度。注意粒子被定义为庞加莱群的表示而不仅是庞加莱代数的表示；换言之，相应的不可约表示必须是 $C\mathcal{PT}$-不变的，其中 $C$ 表示电荷共轭。有质量单粒子态总是 $C\mathcal{PT}$ 不变的，无质量的则不是。例如 $C\mathcal{PT}$ 共轭总是翻转无质量单粒子态螺旋度的符号。出于这个原因，所谓无质量粒子一般指两个螺旋度相反的庞加莱代数无质量表示的组合。由自旋–统计定理，整数自旋/螺旋度的粒子是玻色子，而（严格）半整数自旋/螺旋度的粒子是费米子。

以此类推，在超对称 QFT 的情形，可作如下设定：

% [Tutor Note: 原文此处 $W^\mu=\frac{1}{2}\epsilon^{\mu\nu\rho\sigma}P_\mu M_{\rho\sigma}$ 的指标有重复（$\mu$ 同时出现在 $W^\mu$ 与 $P_\mu$ 上）；按标准定义应为 $W^\mu=\frac{1}{2}\epsilon^{\mu\nu\rho\sigma}P_\nu M_{\rho\sigma}$，这样才能保证 $P_\mu W^\mu=0$ 且 $W^2$ 为卡西米尔算符。此处照原文译出，未擅改。]

\begin{definition}[17.3] 超粒子或超多重态是所论超庞加莱代数的不可约有限维幺正表示，且在 $C\mathcal{PT}$ 共轭下不变。\end{definition}

算符 $P^2$ 在 $4d$ 超庞加莱代数中仍是卡西米尔算符，但 $W^2$ 不是。这可由以下事实看出：超荷 $Q^I_\alpha$ 与 $\bar Q^I_{\dot\alpha}$（它们生成超对称变换）不与洛伦兹群的生成元交换。因此超多重态一般含有质量与内部荷相同、但自旋/螺旋度不同的粒子，因为作用超荷会改变自旋/螺旋度。

超对称 QFT 满足如下重要性质：

\begin{proposition}[17.4] 超对称 QFT 中的每个态都具有非负能量。\end{proposition}

\begin{proof} 取定某个 $I=1,\dots,\mathcal{N}$。有：
\begin{equation}
\begin{aligned}
0\le\sum_\alpha\left\|(Q^I_\alpha)^\dagger|\phi\rangle\right\|^2+\left\|Q^I_\alpha|\phi\rangle\right\|^2
&=\sum_{\alpha=\dot\beta}\langle\phi|\{Q^I_\alpha,\bar Q^I_{\dot\beta}\}|\phi\rangle\\
&=\sum_{\alpha=\dot\beta}2\sigma^\mu_{\alpha\dot\beta}\langle\phi|P_\mu|\phi\rangle\\
&=4\langle\phi|P_0|\phi\rangle=4\langle\phi|H|\phi\rangle\,,
\end{aligned}
\tag{17.2}
\end{equation}
% [原文 p.210]
其中 $H=P_0$ 是哈密顿量。
\end{proof}

在实践中，人们把超对称生成元 $Q^I_\alpha$ 与 $\bar Q^I_{\dot\alpha}$ 重新解释为费米型产生算符 $((a^\dagger_{1,2})^I\propto\bar Q^I_{1,2})$ 与湮灭算符 $(a^I_{1,2}\propto Q^I_{1,2})$，从而得到超庞加莱代数的不可约表示。

\begin{definition}[17.5] 若庞加莱多重态被所有 $a_I$ 湮灭，则称它为一个 Clifford 真空。\end{definition}

对给定的 Clifford 真空，用产生算符 $(a^\dagger)^I$ 以一切可能的方式作用于其上便构造出相应的表示。Clifford 真空生成该表示的最低权空间。

\subsubsection*{无质量超多重态}

有两个重要性质刻画了所有可能超多重态中的无质量超多重态。第一，在无质量粒子四维动量读到 $P_\mu=(E,0,0,E)$ 的参考系中，有：
\[
\{Q^I_\alpha,\bar Q^J_{\dot\beta}\}=\begin{pmatrix}0&0\\ 0&4E\end{pmatrix}_{\alpha\dot\beta}\delta^{IJ}\,.
\]
这意味着只有一半产生算符（即 $(a^\dagger_{1,2})^I$）进入超多重态的构造。换言之，无质量超多重态比一般超多重态更短：它们含有的庞加莱粒子数少于一般预期的数目。每个被非平凡实现的产生算符都把螺旋度提高二分之一，这意味着：若 Clifford 真空的螺旋度为 $j_0$，则相应的超多重态由 $\binom{\mathcal{N}}{k}$ 个螺旋度为 $j_0+k/2$（$k=0,\dots,\mathcal{N}$）的态组成。

无质量超多重态的第二个显著特征与 $C\mathcal{PT}$ 不变性有关。我们已经强调，超庞加莱代数（algebra）的无质量表示通常需要与其螺旋度翻转的对应物配对，才能保持 $C\mathcal{PT}$ 对称性。这意味着，典型地，螺旋度为 $h$ 的无质量粒子必须在同一超多重态内与螺旋度为 $-h$ 的粒子配对。注意当 $j_0=-\mathcal{N}/4$ 时，有可能在不配对的情况下实现 $C\mathcal{PT}$ 不变性；然而这并不总是可行，因为还须满足依赖于超多重态具体结构的其他约束。

后文讨论中将要用到的无质量超多重态属于以下类型，其中粗体数字表示庞加莱成分的螺旋度，而括号前的非粗体数字表示重数：

\begin{itemize}[leftmargin=2em]
\item $\mathcal{N}=1$ 手征多重态：$(\mathbf{0},\mathbf{1/2})\oplus_{C\mathcal{PT}}(\mathbf{0},-\mathbf{1/2})$，
\item $\mathcal{N}=1$ 矢量多重态：$(\mathbf{1/2},\mathbf{1})\oplus_{C\mathcal{PT}}(-\mathbf{1/2},-\mathbf{1})$，
\item $\mathcal{N}=2$ 超多重态（hypermultiplet）：$(-\mathbf{1/2},\mathbf{2}(\mathbf{0}),\mathbf{1/2})\oplus_{C\mathcal{PT}}(\mathbf{1/2},\mathbf{2}(\mathbf{0}),-\mathbf{1/2})$。这里 $j_0=-\mathcal{N}/4$，但为了 $C\mathcal{PT}$ 不变性仍然需要配对，
% [原文 p.211]
\item $\mathcal{N}=2$ 矢量多重态：$(\mathbf{0},\mathbf{2}(\mathbf{1/2}),\mathbf{1})\oplus_{C\mathcal{PT}}(\mathbf{0},\mathbf{2}(-\mathbf{1/2}),-\mathbf{1})$，
\item $\mathcal{N}=4$ 矢量多重态：$(-\mathbf{1},\mathbf{4}(-\mathbf{1/2}),\mathbf{6}(\mathbf{0}),\mathbf{4}(\mathbf{1/2}),\mathbf{1})$。该超多重态是自 $C\mathcal{PT}$ 共轭的。
\end{itemize}

\subsubsection*{有质量超多重态}

相反地，对有质量超多重态，一般没有哪个产生算符被平凡实现，这类超庞加莱代数的表示被称为长的。$(a^\dagger_1)^I$ 降低自旋，而 $(a^\dagger_2)^I$ 提高自旋，增量都是 $1/2$。由此得到 $2^{2\mathcal{N}}$ 个态，其中一半是玻色型、另一半是费米型。在非扩张情形 $\mathcal{N}=1$ 中，两个最重要的有质量超多重态如下（此时括号中的粗体数字表示自旋，前面的非粗体数字表示重数）：

\begin{itemize}[leftmargin=2em]
\item $\mathcal{N}=1$（有质量）物质多重态：$(-\mathbf{1/2},\mathbf{2}(\mathbf{0}),\mathbf{1/2})$，
\item $\mathcal{N}=1$（有质量）矢量多重态：$(-\mathbf{1},\mathbf{2}(-\mathbf{1/2}),\mathbf{2}(\mathbf{0}),\mathbf{2}(\mathbf{1/2}),\mathbf{1})$。该物质内容与一个无质量 $\mathcal{N}=1$ 手征超多重态和一个无质量 $\mathcal{N}=1$ 矢量超多重态的组合相符，正如超希格斯机制所预期的那样。
\end{itemize}

\subsubsection*{BPS 态}

现在考虑扩张超对称（$\mathcal{N}>1$）的情形，并为简单起见、也因为这些正是我们最关心的情形，假定 $\mathcal{N}$ 为偶数。回忆 $\{Q^I_\alpha,Q^J_\beta\}=2\sqrt{2}\epsilon_{\alpha\beta}\mathcal{Z}^{IJ}$，其中 $\mathcal{Z}^{IJ}$ 是代数的复中心荷构成的反对称矩阵。在 $\mathrm{U}(\mathcal{N})$ 中某元素共轭的意义下，可假定 $\mathcal{Z}^{IJ}$ 为分块对角，其 $2\times 2$ 块为 $\Mat(0,\mathcal{Z}_i,-\mathcal{Z}_i,0)$，其中 $\mathcal{Z}_i$ 为实正数，$i=1,\dots,\mathcal{N}/2$。对一切 $i=1,\dots,\mathcal{N}/2$，令
\[
a^i_\alpha=Q^{2i-1}_\alpha+\epsilon_{\alpha\beta}(Q^{2i}_\beta)^\dagger\,,\qquad\text{以及}\qquad
b^i_\alpha=Q^{2i-1}_\alpha-\epsilon_{\alpha\beta}(Q^{2i}_\beta)^\dagger\,.
\]
在质量为 $m$ 的超多重态的静止系中，有
\[
\{a^i_\alpha,(a^j_\beta)^\dagger\}=4(m+\sqrt{2}\mathcal{Z}_i)\delta_{ij}\delta_{\alpha\beta}\,,
\]
\[
\{b^i_\alpha,(b^j_\beta)^\dagger\}=4(m-\sqrt{2}\mathcal{Z}_i)\delta_{ij}\delta_{\alpha\beta}\,,
\]
其余所有反对易子都为零。

\begin{proposition}[17.6] $\mathcal{N}$-扩张超庞加莱代数的任何表示的质量都满足：
\begin{equation}
m\ge\sqrt{2}\mathcal{Z}_i\,,\qquad i=1,\dots,\mathcal{N}/2\,.
\tag{17.3}
\end{equation}
\end{proposition}

\begin{proof} 这可由把式 (17.2) 中的 $Q^I_\alpha$ 换成 $b^i_\alpha$ 得到。更一般地，若如上取 $\mathcal{Z}^{IJ}$ 为分块对角但系数 $\mathcal{Z}_i$ 为复数，
% [原文 p.212]
则通过考虑 $a^i_\alpha$ 与 $b^i_\beta$ 的复系数线性组合便得到不等式 $m\ge\sqrt{2}|\mathcal{Z}_i|$。从现在起，我们把 $\mathcal{Z}_i$ 视为任意复数。
\end{proof}

当 $m$ 对某些 $i$ 等于 $\sqrt{2}|\mathcal{Z}_i|$ 时，某些产生算符被平凡实现。这类非一般的有质量超多重态将比一般的更短。

\begin{definition}[17.7] 若对某些 $i$ 有 $m=\sqrt{2}|\mathcal{Z}_i|$，则称该有质量超多重态为 BPS 态，或短表示。更确切地说，若有 $k$ 个 $i$ 使 $m=\sqrt{2}|\mathcal{Z}_i|$，则称此超多重态为 $k/\mathcal{N}$-BPS 态。反之，若对所有 $i$ 都有 $m>\sqrt{2}\mathcal{Z}_i$，则称该超多重态为长的。\end{definition}

在 $k/\mathcal{N}$-BPS 态中，用超荷构造的 $2\mathcal{N}$ 个费米型产生算符（即当 $\mathcal{Z}_i$ 为实数时的 $(a^i_\alpha)^\dagger$ 与 $(b^i_\alpha)^\dagger$）里有 $2k$ 个（在诸 $(b^i_\alpha)^\dagger$ 之中）被平凡实现。于是相应的表示由 $2^{2(\mathcal{N}-k)}$ 个（庞加莱）态组成，其中一半为玻色型、另一半为费米型。注意 $\frac{1}{2}$-BPS 超多重态所含的态数与无质量超多重态一样多。

带扩张超对称的 $4d$ 超庞加莱代数的有质量表示的例子包括：

\begin{itemize}[leftmargin=2em]
\item $\mathcal{N}=2$ 长矢量多重态：$(-\mathbf{1},\mathbf{4}(-\mathbf{1/2}),\mathbf{6}(\mathbf{0}),\mathbf{4}(\mathbf{1/2}),\mathbf{1})$。一个 Clifford 真空自旋不为零的 $\mathcal{N}=2$ 长多重态必然含有自旋严格大于 1 的粒子。
\item $\mathcal{N}=2$ 半 BPS 有质量超多重态：$(\mathbf{2}(-\mathbf{1/2}),\mathbf{4}(\mathbf{0}),\mathbf{2}(\mathbf{1/2}))$。态的加倍由 $\SU(2)_R$ 不变性所要求：为保持 $\SU(2)_R$，超多重态中的标量必须按 $\SU(2)_R$ 的双重态（四元数）表示变换。然而，由产生–湮灭代数作用在 Clifford 真空上所得的两个标量是实的，因而无法满足这一要求（第 18 章将对此作更多讨论）。这也解释了上文提到的无质量 $\mathcal{N}=2$ 超多重态中态的加倍，因为 $C\mathcal{PT}$ 自共轭会要求两个标量为实。
\item $\mathcal{N}=2$ 半 BPS 矢量多重态：$(-\mathbf{1},\mathbf{2}(-\mathbf{1/2}),\mathbf{2}(\mathbf{0}),\mathbf{2}(\mathbf{1/2}),\mathbf{1})$。
\item $\mathcal{N}=4$ 四分之一 BPS 矢量多重态：$(-\mathbf{1},\mathbf{4}(-\mathbf{1/2}),\mathbf{6}(\mathbf{0}),\mathbf{4}(\mathbf{1/2}),\mathbf{1})$。
\end{itemize}

\subsection{共形代数与超共形代数}

\subsubsection*{共形不变性}

在平直空间中推广庞加莱对称性的另一条途径是考虑标度不变理论。标度不变性与 $S$-矩阵的定义不相容，因此这类理论不受 Coleman–Mandula 定理的约束。

% [原文 p.213]
一大类标度不变理论是共形场论（CFT），其玻色型对称群包含闵可夫斯基时空 $\mathcal{M}_d$ 的共形群。反过来，我们所知的多数标度不变理论实际上都是 CFT。更多细节参见 [164, 108, 268, 6]。

\begin{definition}[17.8] $\mathcal{M}_d$ 的共形群是使度量在相差一个标度因子的意义下不变的变换群：$\eta_{\mu\nu}\to\Omega^2(x)\eta_{\mu\nu}$。它包含庞加莱变换、伸缩 $x^\mu\to\lambda x^\mu$（$\lambda\in\R_{>0}$）、反演 $x^\mu\to x^\mu/x^2$，以及由平移与反演共轭得到的特殊共形变换。在无穷小对称性的层面，共形代数 $\so(d,2)$ 由无穷小平移 $P_\mu$、洛伦兹变换 $M_{\mu\nu}$、伸缩 $D$ 与特殊共形变换 $K_\mu$ 生成，它们特别满足：
\[
[D,P_\mu]=-\ii P_\mu\,,\qquad [D,K_\mu]=\ii K_\mu\,.
\]
\end{definition}

在共形代数中，$P^2$ 不是卡西米尔算符。作为伸缩算符本征函数的场，即满足 $\phi(x)\to\lambda^\Delta\phi(\lambda x)$ 的场，称为维数为 $\Delta$ 的准初级场；而同时被诸 $K_\mu$ 湮灭的准初级场称为初级场。共形代数的表示通过庞加莱生成元作用在初级场（“最高权”矢量）上来构造。

在幺正 CFT 中，场的维数有下界，且依赖于它所在的庞加莱代数表示。例如在一般维数 $d$ 中：
\[
\Delta(\text{标量 }\phi)\ge(d-2)/2\,,
\]
\[
\Delta(\text{旋子 }\psi)\ge(d-1)/2\,,
\]
\[
\Delta(\text{矢量 }J_\mu)\ge d\,,
\]
等号分别当 $\partial^2\phi=0$（自由标量）、$\partial\!\!\!/\psi=0$（自由费米场）与 $\partial^\mu J_\mu=0$（守恒流）时成立。这些不等式称为幺正性下界。

\subsubsection*{超共形代数}

把超对称与共形不变性结合起来便得到超共形代数 [164, 268, 6]。更确切地说，$d$ 维超共形代数须同时包含共形代数 $\so(d,2)$ 与带 $\mathcal{N}$ 个超对称的 $d$ 维庞加莱超代数，且其费米型生成元须处于 $\so(d,2)$ 的旋子表示中。

例如当 $d=4$ 时，在共形生成元 $P_\mu$、$M_{\mu\nu}$、$D$、$K_\mu$ 与超对称生成元 $Q_\alpha$、$\bar Q_{\dot\alpha}$ 之上，超共形代数还容许超共形生成元 $S_\alpha$、$\bar S_{\dot\alpha}$ 以及 $Q$ 与 $S$ 的换位子中的 R 对称性生成元，其中 R 对称性在 17.1.3 节定义。

一般地，超共形代数只存在于维数 $d\le 6$ 的情形，且当 $d\ge 3$ 时受到很强的约束。对 $d=4$，带 $\mathcal{N}$ 个超对称的超共形代数只有：

% [原文 p.214]
\begin{itemize}[leftmargin=2em]
\item $\mathcal{N}\le 3$：$\su(2,2|\mathcal{N})\supset\su(2,2)\times\mathfrak{u}(\mathcal{N})_R$，
\item $\mathcal{N}=4$：$\mathfrak{psu}(2,2|4)\supset\su(2,2)\times\su(4)_R$。
\end{itemize}

在 SCFT 中，庞加莱超对称的数目同样有上界；在 $d=4$ 时最大数目为 $\mathcal{N}=4$，对应 32 个超荷（8 个 $Q^I_\alpha$、8 个 $\bar Q^I_{\dot\alpha}$、8 个 $S^I_\alpha$ 与 8 个 $\bar S^I_{\dot\alpha}$）。超共形初级场是被所有 $K_\mu$ 与超共形生成元湮灭的准初级场；手征初级场是被所有 $\bar Q_{\dot\alpha}$ 湮灭的初级场。

\subsection{内部对称性}

量子场论中的内部对称性群囊括了一切非时空对称性。在这个群中，可以辨认出 R 对称性，其定义是那些不与超荷交换的对称性。相反地，所有非 R 对称性都与超庞加莱代数交换。

内部对称性群刻画一个量子场论，并且无论理论的具体表述如何都保持不变。这与规范对称性形成对照：后者是理论描述的特征，而不是理论本身的性质。例如，同一个量子场论可以有不同的规范理论对偶描述，其规范群各不相同。我们将在本章末尾讨论 Seiberg 对偶时遇到这一现象的第一个例子，后文还会遇到更多例子。然而，同一理论的对偶描述必定具有相同的内部对称性。

\subsubsection*{整体对称性与局域对称性}

内部对称性常被称为整体的，因为它们同时作用于所有时空点。例如，质量为 $m$ 的自由复标量场 $\phi(x)$ 的 $4d$ 理论由作用量
\begin{equation}
S=\int \dd^4x\,\partial_\mu\bar\phi\,\partial^\mu\phi+m^2|\phi|^2\,,
\tag{17.4}
\end{equation}
定义，它具有由 $\phi(x)\to\ee^{\ii\alpha}\phi(x)$ 给出的 $\mathrm{U}(1)$ 整体对称性，其中 $\alpha$ 是常数。

可以通过允许对称性非平凡地依赖于时空来“规范化”它，即取 $\phi(x)\to\ee^{\ii\alpha(x)}\phi(x)$，其中 $\alpha(x)$ 是 $x$ 的函数。这样的时空依赖对称性称为局域的。规范化（gauging）需要引入一个 $\mathrm{U}(1)$ 矢量场（或规范场）$A_\mu(x)$ 来抵消由此产生的额外项。若在作用量中把导数 $\partial_\mu$ 换成协变导数 $D_\mu=\partial_\mu+\ii A_\mu$，并且当 $\phi(x)\to\ee^{\ii\alpha(x)}\phi(x)$ 时 $A_\mu$ 按 $A_\mu\to A_\mu-\partial_\mu\alpha$ 变换，则 $D_\mu\phi\to\ee^{\ii\alpha(x)}D_\mu\phi$，且修改后的作用量
\[
S=\int \dd^4x\,D_\mu\bar\phi\,D^\mu\phi+m^2|\phi|^2
\]
在局域 $\mathrm{U}(1)$ 对称性下不变。再加上矢量场 $A_\mu$ 的标准动能项 $\propto F_{\mu\nu}F^{\mu\nu}$（其中 $F_{\mu\nu}=\partial_\mu A_\nu-\partial_\nu A_\mu$），便得到标量电动力学的作用量。从几何上看，$A_\mu$ 是一个联络 1-形式。人们区分小与大的规范变换：前者与恒等变换连续相连，后者则不然。

% [Tutor Note: 原文此处 (17.4) 与紧随其后的修改作用量都把 $m^2|\phi|^2$ 写在积分号之后（排版上属于被积函数）。核对质量量纲可验证：$[\phi]=1$、$[\partial_\mu]=1$，故 $m^2|\phi|^2$ 是 4 维密度的一部分，应理解为 $S=\int\dd^4x\,(\partial_\mu\bar\phi\partial^\mu\phi+m^2|\phi|^2)$。译文照原文排版，未改公式。]

% [原文 p.215]
与整体对称性相反，规范对称性并非真正的对称性：系统的所有态与可观测量都必须是规范不变的。把整体对称性规范化，本质上相当于用该对称性对态与可观测量作商，从而得到不再具有原对称性的新理论。尽管如此，规范对称性确实有物理后果。例如，规范对称性诱导出类高斯定律，并且是构造拉格朗日式渐近自由理论的必要条件，这一结果称为 Coleman–Gross 定理 [95]。

\begin{remark}[17.9] 规范对称性频繁出现在含微分形式 $A^{(k)}$（次数 $k\ge 1$）的 $d$ 维量子场论中，这些形式通过其外微分 $\dd A^{(k)}$ 进入作用量。由于性质 $\dd\circ\dd=0$，作用量在变换 $A^{(k)}\to A^{(k)}+\dd B^{(k-1)}$ 下保持不变，其中 $B^{(k-1)}$ 是任意 $(k-1)$ 次微分形式。这个对称性是局域的而非整体的，因为 $B^{(k-1)}$ 不必是闭的（这是“函数为常数”在更高次形式中的类比）。当 $k>1$ 时，这称为高阶形式规范对称性。\end{remark}

\subsubsection*{作为拓扑算符的整体对称性}

对整体对称性的历史理解深深植根于诺特定理的第一定理。按该定理，每个连续整体对称性群 $G$ 都给出一个守恒流：一个满足守恒方程 $\partial^\mu j^{(1)}_\mu(x)=0$ 的 1-形式 $j^{(1)}_\mu(x)$（参见例如 [108, Chap. 2]）。例如在由式 (17.4) 定义的理论中，前文讨论的整体 $\mathrm{U}(1)$ 对称性所对应的守恒流取下形式：
\[
j^{(1)}_\mu(x)=\phi(x)\partial_\mu\bar\phi(x)-\bar\phi(x)\partial_\mu\phi(x)\,.
\]

流 $j^{(1)}_\mu(x)$ 的守恒方程可用微分形式的语言更自然地改写：$\dd(\star_d j^{(1)})=0$，其中 $\star_d$ 是 $d$ 维 Hodge 星算子。因此 $\star_d j^{(1)}$ 是一个 $(d-1)$-形式。它可在余维 1 子流形 $M_{(d-1)}$ 上积分，得到相应的守恒荷：
\[
Q(M_{(d-1)})=\int_{M_{(d-1)}}\star_d j^{(1)}\,.
\]
守恒方程意味着：只要 $M_{(d-1)}$ 与 $N_{(d-1)}$ 能连续形变为彼此，就有 $Q(M_{(d-1)})=Q(N_{(d-1)})$；荷 $Q$ 的守恒正是在这个意义上说的。

% [原文 p.216]
若所论对称性对应于李群 $G$，则流 $j^{(1)}_\mu$ 实际上取值于李代数 $\mathfrak{g}=\mathrm{Lie}(G)$，$G$ 的每个生成元对应一个标量守恒荷。通过指数化，可以构造与子流形 $M_{(d-1)}$ 及任意 $g\in G$ 相联系的拓扑算符或缺陷 $\mathcal{T}_g(M_{(d-1)})$。$G$ 对称性的带荷算符是由 $G$ 的表示 $R$ 刻画的场 $\mathcal{O}^i(x)$（它们是局域的，即时空上的函数）。该对称性转化为关联函数之间的 Ward–Takahashi 恒等式：
\begin{equation}
\langle\mathcal{T}_g(S_{(d-1)})\mathcal{O}^i(x)\rangle=\left(\Mat_R(g)\right)^i_{\ j}\langle\mathcal{O}^j(x)\rangle\,,
\tag{17.5}
\end{equation}
其中 $\Mat_R(g)$ 是 $g$ 在表示 $R$ 中的像，$S_{(d-1)}$ 是包围点 $x$ 的拓扑 $(d-1)$-维球面。支撑在同伦等价子流形上的拓扑缺陷可按某种方式融合，该方式对应于 $G$ 中的乘法。事实上，人们可以把通常的对称性定义如下：

\begin{definition}[17.10] 量子场论中对应于群 $G$ 的对称性是一族拓扑缺陷 $\mathcal{T}_g([M_{(d-1)}])$，以 $G$ 的元素 $g$ 为指标，其中 $[M_{(d-1)}]$ 是时空余维 1 子流形的任一同伦类。\end{definition}

尽管上述构造依赖于守恒流的存在（而守恒流只在 $G$ 连续时才存在），对应于整体对称性的拓扑缺陷在 $G$ 离散时也存在，并在时空余维 1 子流形的同伦类上实现 $G$ 的元素。$G$ 中的群乘法反映在这些拓扑算符的融合规则中。

\subsubsection*{广义整体对称性}

对应于通常整体对称性的拓扑缺陷——本节余下部分我们称之为 0-形式对称性——支撑在余维 1 流形上并作用于局域算符。可以放松对拓扑缺陷余维数的限制，允许它们支撑在余维 $(p+1)$ 流形上，其中 $0\le p\le d-1$。这一观点最先在 [149] 中倡导，启发了 QFT 中标准对称性概念的富有成果的推广。这里我们简要回顾这些推广，详细的论述可参见例如 [53, 327] 及其中的参考文献。

\begin{definition}[17.11] 这类算符编码所谓的 $p$-形式对称性；它们作用于延展的 $p$ 维对象，是式 (17.5) 的推广。\end{definition}

\begin{remark}[17.12] 一个版本的 Eckmann–Hilton 论证表明，$p\ge 1$ 的 $p$-形式对称性总是交换的：实现该对称性的拓扑缺陷集合所满足的融合规则对应于某个交换群 $G$ 中的乘法。\end{remark}

连续的 $p$-形式对称性对应于守恒的诺特 $(p+1)$-形式流 $j^{(p+1)}$，满足 $\dd\star_d j^{(p+1)}=0$；在余维 $(p+1)$ 流形上积分 $\star_d j^{(p+1)}$ 便定义拓扑荷，而指数化之后，拓扑

% [原文 p.217]
缺陷。如前所述，守恒流的存在并不等同于对称性，后者完全可能是离散的。

让我们考虑 $d=4$ 的纯 $\mathrm{U}(1)$ 规范理论（即麦克斯韦理论），其作用量为
\[
S[A^{(1)}]\propto\int \dd^4x\,F^{(2)}\wedge\star_4 F^{(2)}=\frac{1}{2}\int \dd^4x\,F_{\mu\nu}F^{\mu\nu}\,.
\]
运动方程 $\dd\star_4 F^{(2)}=0$ 可视为一个 2-形式流的守恒方程；相应的 $\mathrm{U}(1)_e^{(1)}$ 1-形式对称性称为电的，它编码了把规范场 $A^{(1)}$ 平移一个闭 1-形式 $\lambda^{(1)}$（即 $\dd\lambda^{(1)}=0$）的自由度。如上文所述，“闭”是“函数为常数”在更高次微分形式中的类比；在这个意义上，$\mathrm{U}(1)_e^{(1)}$ 是整体对称性。2-形式 $\star_4 F^{(2)}$ 可在余维 2 流形上积分，从而得到守恒荷（高斯定律）与余维 2 拓扑缺陷。

这可推广到 $d$ 维纯 $\mathrm{U}(1)$ 高阶规范理论，其作用量定义为
\[
S[A^{(k)}]\propto\int \dd^dx\,F^{(k+1)}\wedge\star_d F^{(k+1)}\,,
\]
其中 $F^{(k+1)}=\dd A^{(k)}$：这类理论容许一个电的 $k$-形式对称性 $\mathrm{U}(1)_e^{(k)}$。

这类理论还显示出另一个拓扑性质的守恒方程：
\[
\dd F^{(k+1)}\propto\dd\star_d\left(\star_d F^{(k+1)}\right)=0\,.
\]
守恒流是一个 $(d-k-1)$-形式，而对称性是 $\mathrm{U}(1)_m^{(d-k-2)}$，称为磁的。它对应于把磁对偶规范场 $A_D^{(d-k-2)}$ 平移一个闭 $(d-k-2)$-形式的自由度。我们将在第 19 章更详细地讨论电磁对偶。

在 $4d$ 纯 $\mathrm{U}(1)$ 规范理论中，实现 $\mathrm{U}(1)_e^{(1)}$ 电 1-形式对称性的拓扑缺陷支撑在余维 2 流形上。带荷对象是 Wilson 算符 $W_L(q)=\exp(\ii q\oint_L A^{(1)})$，其中 $L$ 是时空的 1 维子流形，$q$ 是电荷。当 $A^{(1)}\to A^{(1)}+\lambda^{(1)}$ 且 $\lambda^{(1)}$ 闭时，有 $W_L(q)\to\ee^{\ii qM(\lambda^{(1)},L)}W_L(q)$，其中 $M(\lambda^{(1)},L)$ 是 $\lambda^{(1)}$ 沿 $L$ 的单值性。同样地，磁 1-形式对称性的带荷对象是 't Hooft 算符。

当存在带电荷和/或带磁荷的场时，1-形式对称性 $\mathrm{U}(1)_e^{(1)}\times\mathrm{U}(1)_m^{(1)}$ 被屏蔽破缺为子群。例如，当存在电荷为 $n$ 的带电场 $\phi$ 时，由于电荷为 $n$ 的 Wilson 线可以终止于 $\phi(x)$，电 1-形式对称性破缺为其离散子群 $(\Z_n)_e^{(1)}$。图 17.1 给出了这一情形的图示。

\begin{remark}[17.13] 四维非阿贝尔杨–米尔斯理论同样容许电与磁的 1-形式对称性，其带荷对象是 Wilson–'t Hooft 线算符。它们编码了关于规范群整体结构的精细信息，而纯杨–米尔斯作用量

% [原文 p.218]
\Fig[0.86\linewidth]{figures/fig_17_1.png}{17.1}{带荷算符存在时 1-形式对称性的破缺。}

\[
S=\frac{-1}{2g^2}\int \dd^4x\,\tr_{\mathfrak{g}}(F\wedge F)\,,
\]
对此并不敏感，因为它只依赖于李代数 $\mathfrak{g}$。例如 $\SU(N)$ 与 $\mathrm{PSU}(N)$ 杨–米尔斯理论共享同一作用量：它们的局域动力学相同。与代数 $\mathfrak{g}$ 相对应的规范群 $G$ 的整体结构之选取，等价于局部 Wilson–'t Hooft 算符的一个极大集合之选取 [7]，而这又决定了理论的 1-形式整体对称性 [149]。\end{remark}

近年来，量子场论与规范理论的广义对称性得到了细致研究，并从中涌现出有趣的数学结构。不同阶的整体对称性可以组合，导致（范畴意义上的）高群；此外，许多规范理论容许不可逆拓扑缺陷，它们推广了二维 Ising 模型的 Kramers–Wannier 对偶缺陷。

\subsubsection*{背景场与规范化}

让我们考虑由作用量 $S$ 描述、具有连续 $p$-形式整体对称性 $G^{(p)}$ 及相应诺特流 $j^{(p+1)}$ 的 $d$ 维量子场论，后者由变换规律
\begin{equation}
S\to S+\ii\int \dd\lambda^{(p)}\wedge\star_d j^{(p+1)}\,,
\tag{17.6}
\end{equation}
定义，其中 $\lambda^{(p)}$ 是一般的（时空依赖的）$p$-形式对称性参数。

可以考虑理论处于整体对称性 $G^{(p)}$ 的背景场 $A^{(p+1)}$ 之中，即考虑修改后的作用量：
\begin{equation}
S[A^{(p+1)}]=S-\ii\int A^{(p+1)}\wedge\star_d j^{(p+1)}\,.
\tag{17.7}
\end{equation}
在背景场存在时，若要求 $A^{(p+1)}$ 按 $A^{(p+1)}\to A^{(p+1)}+\dd\lambda^{(p)}$ 变换，则对任意（不必闭的）$p$-形式 $\lambda^{(p)}$，理论不变。

% [Tutor Note: 式 (17.6)–(17.7) 的量纲与次数核对：$j^{(p+1)}$ 是 $(p+1)$-形式，$\star_d$ 把它变成 $(d-p-1)$-形式；于是 $\dd\lambda^{(p)}\wedge\star_d j^{(p+1)}$ 与 $A^{(p+1)}\wedge\star_d j^{(p+1)}$ 都是 $d$-形式，可在 $d$ 维时空上积分。注意 $\lambda^{(p)}$ 只须是 $p$-形式，不必闭；闭性仅在要求 $A$ 的平移不改变场强时才需要。]
% [原文 p.219]
但要注意 $A^{(p+1)}$ 是修改后理论的参数而非动力学场：人们不对 $A^{(p+1)}$ 求和，因此在固定背景 $A^{(p+1)}$ 的修改后理论中，局域对称性并不是该理论的对称性。

\begin{definition}[17.14] 对所有可能的背景 $A^{(p+1)}$ 求和，是得到另一个理论——其中局域对称性真正成为规范对称性——的途径；这是描述规范化程序的途径之一。\end{definition}

对应于 $G^{(p)}$ 的拓扑缺陷写为：
\[
\begin{aligned}
\mathcal{T}_\alpha(M_{(d-p-1)})&=\exp\left\{\ii\alpha\int_{M_{(d-p-1)}}\star_d j^{(p+1)}\right\}\\
&=\exp\left\{\ii\alpha\int\Pi^{(d-p-1)}_M\wedge\star_d j^{(p+1)}\right\},
\end{aligned}
\]
其中 $\Pi^{(d-p-1)}_M$ 是与 $M_{(d-p-1)}$ 庞加莱对偶的微分形式。这表明，插入实现 $G^{(p)}$ 的拓扑缺陷对应于开启一个背景规范场。因此，规范化程序可以等价地描述为对时空中一切可能的对称性缺陷网络求和。

\begin{definition}[17.15] 规范化存在障碍：'t Hooft 反常。具有守恒流 $j^{(p+1)}$ 的广义对称性有 't Hooft 反常，如果 $\dd\star j^{(p+1)}$ 以本质的方式依赖于背景场。由于后者是非动力学的，'t Hooft 反常是整体对称性的特征。离散对称性也有 't Hooft 反常，并且可能涉及两个或多个不同的整体对称性，此时人们称之为混合 't Hooft 反常。\end{definition}

't Hooft 反常对量子场论的动力学给出严格约束，因为它们在重整化群流下守恒：这一结果称为 't Hooft 反常匹配条件。一个著名的推论是：若某个通常的整体对称性在量子场论中有非零 't Hooft 反常，则该理论不可能有质量间隙（意思是 QFT 中任一态的能量与真空能量之差下方被一个严格正数所界），因为那样就无法在低能有效理论中匹配该反常。't Hooft 反常是量子场论的内禀性质，而其他类型的反常（例如手征整体反常或规范反常）则更像是某一具体理论描述的特征。我们将在 17.3.3 节回到对反常的更详细讨论。

\section{经典场论}

本节我们探讨 $4d$ $\mathcal{N}=1$ 规范理论的拉格朗日密度的构造，关注其经典结构。这些理论具有若干关键特征，例如存在超对称真空的经典模空间。

% [原文 p.220]
经典的 $4d$ $\mathcal{N}=1$ 规范理论常常通向同一理论的良好定义的量子版本。虽然并非所有量子场论（QFT）都可用拉格朗日量描述，但当拉格朗日量存在时，它就是研究该理论的有力工具。特别地，我们强调可重整化的 QFT，因为它们在后面将具有首要的重要性。相应地，我们将区分两类拉格朗日密度：可重整化的与不可重整化的。虽然可重整化是一个量子概念，但当拉格朗日量存在时，它可由拉格朗日结构推断出来。可重整化的拉格朗日量描述紫外（微观）理论，而不可重整化的拉格朗日量通常编码只在某一能标之下才有效的低能有效理论。

在本节及下一节中，我们聚焦于 $4d$ $\mathcal{N}=1$ 理论，主要遵循 [50]。

\subsection{超场}

把狭义相对论与量子力学结合起来所施加的约束，使得我们有必要考虑携带对称代数表示的场，而不仅仅是关注前文 17.1.1 节讨论过的单粒子表示。场提供了更灵活的框架，因为它们描述多粒子系统并能说明产生与湮灭过程，同时尊重时空与超对称的底层对称性（参见例如 [19]）。庞加莱代数的生成元在场空间中被表示为微分算符，例如 $P_\mu\leftrightarrow-\ii\partial_\mu$：在一阶近似下，无穷小平移 $x\to x+a$ 下，场 $\phi(x)$ 的变化为
\[
\phi(x)\to\phi(x+a)=\phi(x)+a^\mu\partial_\mu\phi(x)=\phi(x)+\ii a^\mu P_\mu\phi(x)\,.
\]
类似地，在超对称（SUSY）相对论性量子理论中，人们关心超庞加莱代数在场上的表示。17.1.1 节的超多重态联系于超场，后者（在量子化之后）描述相应超多重态的产生与湮灭。

超场与超对称拉格朗日密度的构造在 $4d$ $\mathcal{N}=1$ 超空间 $\mathcal{M}_{4|1}$ 中大大简化。设 $x^\mu,\theta_\alpha,\bar\theta_{\dot\alpha}$ 是 $\mathcal{M}_{4|1}$ 上的标准坐标。超荷 $Q_\alpha$ 与 $\bar Q_{\dot\alpha}$ 在超空间中体现为微分算符：
\[
Q_\alpha=-\ii\partial_\alpha-\sigma^\mu_{\alpha\dot\beta}\bar\theta^{\dot\beta}\partial_\mu\,,\qquad
\bar Q_{\dot\alpha}=\ii\bar\partial_{\dot\alpha}+\theta^\beta\sigma^\mu_{\beta\dot\alpha}\partial_\mu\,,
\]
\[
\Rightarrow\quad \{Q_\alpha,\bar Q_{\dot\beta}\}=-2\ii\sigma^\mu_{\alpha\dot\beta}\partial_\mu=2\sigma^\mu_{\alpha\dot\beta}P_\mu\,.
\]

\begin{definition}[17.16] 超场是 $\mathcal{M}_{4|1}$ 上的函数 $Y(x,\theta,\bar\theta)$，它在旋子参数 $\epsilon_\alpha,\bar\epsilon_{\dot\alpha}$ 的无穷小超平移下按如下方式变换：

% [Tutor Note: 上式 $Q_\alpha$ 与 $\bar Q_{\dot\alpha}$ 的符号约定决定了 (17.9) 中各分量项的相对符号，核对时可用 $\{Q_\alpha,\bar Q_{\dot\beta}\}=2\sigma^\mu_{\alpha\dot\beta}P_\mu$ 反推：把 $P_\mu=-\ii\partial_\mu$ 与 $\bar Q$ 的显式表达式代入反对易子，应重新得到同一结果。]
% [原文 p.221]
\[
Y(x,\theta,\bar\theta)\to Y(x+\delta x,\theta+\delta\theta,\bar\theta+\delta\bar\theta)=(\ii\epsilon Q+\ii\bar\epsilon\bar Q)Y(x,\theta,\bar\theta)\,,
\]
其中 $\delta\theta=\epsilon$，$\delta\bar\theta=\bar\epsilon$，$\delta x^\mu=\ii\theta\sigma^\mu\bar\epsilon-\ii\epsilon\sigma^\mu\bar\theta$，这由式 (17.1) 中的关系 $\{Q_\alpha,\bar Q_{\dot\beta}\}\propto P_\mu$ 所要求。\end{definition}

由于 $\theta_\alpha$ 与 $\bar\theta_{\dot\alpha}$ 是 Grassmann 变量，任何超场都可展开为有限多个场分量：
\[
\begin{aligned}
Y(x,\theta,\bar\theta)={}&f(x)+\theta\psi(x)+\bar\theta\bar\chi(x)+\theta\theta m(x)+\bar\theta\bar\theta n(x)+\theta\sigma^\mu\bar\theta v_\mu(x)\\
&+\theta\theta\bar\theta\bar\lambda(x)+\bar\theta\bar\theta\theta\rho(x)+\theta\theta\bar\theta\bar\theta d(x)\,.
\end{aligned}
\]
超场的乘积与和仍是超场。

对应于 $4d$ $\mathcal{N}=1$ SUSY 代数特定不可约表示的超场，特别是手征与矢量超多重态，是通过施加与 $Q_\alpha$ 和 $\bar Q_{\dot\alpha}$ 交换的约束得到的。

\subsubsection*{手征超场}

考虑微分算符
\begin{equation}
D_\alpha=\partial_\alpha+\ii\sigma^\mu_{\alpha\dot\beta}\bar\theta^{\dot\beta}\partial_\mu\,,\qquad
\bar D_{\dot\alpha}=\bar\partial_{\dot\alpha}+\ii\theta^\beta\sigma^\mu_{\beta\dot\alpha}\partial_\mu\,.
\tag{17.8}
\end{equation}
它们满足如下反对易关系：
\[
\{D_\alpha,\bar D_{\dot\beta}\}=-2\sigma^\mu_{\alpha\dot\beta}P_\mu\,,\qquad
\{D_\alpha,D_\beta\}=\{D_\alpha,Q_\beta\}=\{D_\alpha,\bar Q_\beta\}=0\,,
\]
\[
\{\bar D_{\dot\alpha},\bar D_{\dot\beta}\}=\{\bar D_{\dot\alpha},Q_\beta\}=\{\bar D_{\dot\alpha},\bar Q_\beta\}=0\,.
\]

\begin{definition}[17.17] 手征超场是满足对每个 $\dot\alpha$ 都有 $\bar D_{\dot\alpha}\Phi=0$ 的超场 $\Phi$。反手征超场是满足对每个 $\alpha$ 都有 $D_\alpha\Psi=0$ 的超场 $\Psi$。\end{definition}

在坐标中，手征超场写为
\begin{equation}
\begin{aligned}
\Phi(x,\theta,\bar\theta)={}&\phi(x)+\sqrt{2}\theta\psi(x)+\ii\theta\sigma^\mu\bar\theta\partial_\mu\phi(x)-\theta\theta F(x)\\
&-\frac{\ii}{\sqrt{2}}\theta\theta\partial_\mu\psi(x)\sigma^\mu\bar\theta-\frac{1}{4}\theta\theta\bar\theta\bar\theta\,\Box\phi(x)\,,
\end{aligned}
\tag{17.9}
\end{equation}
它对应于 $4d$ $\mathcal{N}=1$ 超庞加莱代数的手征超多重态。用 $(\ii\epsilon Q+\ii\bar\epsilon\bar Q)$ 作用，可推出如下 SUSY 变换：
\[
\delta\phi=\sqrt{2}\epsilon\psi\,,
\]
\[
\delta\psi_\alpha=\sqrt{2}\ii(\sigma^\mu\bar\epsilon)_\alpha\partial_\mu\phi-\sqrt{2}\epsilon_\alpha F\,,
\]
\[
\delta F=\ii\sqrt{2}\partial_\mu\psi\sigma^\mu\bar\epsilon\,.
\]

% [原文 p.222]
手征超场的乘积与和仍是手征超场。此外，若 $\Phi$ 是手征超场，则其共轭 $\bar\Phi$ 是反手征超场。反手征超场具有与式 (17.9) 类似的分量展开，由此可以像对手征超场那样推出超对称变换。

\subsubsection*{矢量超场}

\begin{definition}[17.18] $\mathcal{N}=1$ 矢量超场是满足实性条件 $\bar V=V$ 的超场 $V$。\end{definition}

在分量中可以直接验证出现项 $\theta\sigma^\mu\bar\theta v_\mu$，其中 $v_\mu$ 是实矢量场。此外，若 $\Phi$ 是手征超场，则 $\Phi+\bar\Phi$ 是矢量超场，其相应的矢量场为 $-2\partial_\mu\operatorname{Im}\phi$。因此，出现在 $V+\Phi+\bar\Phi$ 的分量分解中的实矢量场是 $v_\mu-2\partial_\mu\operatorname{Im}\phi$：变换 $V\to V+\Phi+\bar\Phi$ 是阿贝尔规范变换的超对称推广。

\begin{remark}[17.19] 总存在手征超场 $\Phi$，使得 $V\to V_{\mathrm{WZ}}=V+\Phi+\bar\Phi$ 处于称为 Wess–Zumino 规范的正规形式，此时它分解为：
\[
V_{\mathrm{WZ}}(x,\theta,\bar\theta)=\theta\sigma^\mu\bar\theta v_\mu(x)+\ii\theta\theta\bar\theta\bar\lambda(x)-\ii\bar\theta\bar\theta\theta\lambda(x)+\frac{1}{2}\theta\theta\bar\theta\bar\theta D(x)\,.
\]
$V$ 的矢量场 $v_\mu(x)$ 是通常的规范场，马约拉纳费米子 $\lambda$ 称为与 $v_\mu$ 相联系的规范微子，而 $D(x)$ 是 $\mathcal{N}=1$ 矢量超场的辅助场。超场 $V$ 体现了 $4d$ $\mathcal{N}=1$ 超庞加莱代数的矢量超多重态。\end{remark}

\subsubsection*{（超）流超场}

对应于整体连续（非 R）对称性的诺特流 $j_\mu$ 是矢量超场 $\mathcal{J}$ 的一个分量，$\mathcal{J}$ 满足 $D^2\mathcal{J}=\bar D^2\mathcal{J}=0$；这一附加约束把守恒律 $\partial^\mu j_\mu=0$ 超对称化。

超对称还定义可被超对称化的守恒流 $S_{\alpha\mu}$。Ferrara–Zumino 完备化由一对矢量超场 $\mathcal{J}_\mu$ 与手征超场 $X$ 组成，满足 $2\bar D^{\dot\alpha}\sigma^\mu_{\alpha\dot\alpha}\mathcal{J}_\mu=D_\alpha X$。

最后，R 对称性的守恒流被超对称化为矢量超场 $\mathcal{R}_\mu$ 与手征超场 $\chi_\alpha$ 组成的一对，满足 $2\bar D^{\dot\alpha}\sigma^\mu_{\alpha\dot\alpha}\mathcal{R}_\mu=\chi_\alpha$，以及 $\bar D_{\dot\alpha}\bar\chi^{\dot\alpha}-D^\alpha\chi_\alpha=0$。

% [原文 p.223]
\subsection{拉格朗日量}

通过把超场 $\mathcal{A}(x,\theta,\bar\theta)$ 在超空间上积分，可以容易地得到 $\mathcal{N}=1$ SUSY 不变的 $4d$ 作用量：
\[
S=\int \dd^4x\,\dd^2\theta\,\dd^2\bar\theta\,\mathcal{A}(x,\theta,\bar\theta)=\int \dd^4x\,\mathcal{L}(\phi(x),\psi(x),\lambda(x),A_\mu(x),\dots)\,.
\]
作用量的不变性来自测度的不变性，因为
\[
\delta_{\epsilon,\bar\epsilon}S=\int \dd^4x\,\dd^2\theta\,\dd^2\bar\theta\,\delta_{\epsilon,\bar\epsilon}\mathcal{A}(x,\theta,\bar\theta)=0\,,
\]
其中第二个等号来自
\[
\int \dd^2\theta\,\dd^2\bar\theta\,\delta_{\epsilon,\bar\epsilon}\mathcal{A}(x,\theta,\bar\theta)
\]
是全导数。物理上，为使拉格朗日密度有意义，$\mathcal{A}$ 还须满足附加约束。注意可以改写
\[
\int \dd^4x\,\dd^2\theta\,\dd^2\bar\theta\,\mathcal{A}(x,\theta,\bar\theta)=\int \dd^4x\,\dd^2\theta\,\bar D^2\mathcal{A}(x,\theta,\bar\theta)\,,
\]
而 $\bar D^2\mathcal{A}$ 是手征超场，因为 $\bar D^3=0$：一般地，用 $\dd^2\theta$ 积分手征超场也能得到 $\mathcal{N}=1$ 超对称不变的拉格朗日密度。然而，并非所有形如
\[
\int \dd^4x\,\dd^2\theta\,\Phi
\]
的项（其中 $\Phi$ 是手征超场）都能改写为在完整超空间上的积分。

\begin{definition}[17.20] 不能表示为在整个超空间上积分的积分项称为 F 项，而能用这种方式写出的称为 D 项。\end{definition}

\subsubsection*{手征超场}

\begin{definition}[17.21] 一组手征超场 $\Phi_i$ 的动能项是如下形式的 D 项：
\begin{equation}
\int \dd^2\theta\,\dd^2\bar\theta\,K(\Phi_i,\bar\Phi_i)\,,
\tag{17.10}
\end{equation}
其中 $K$ 是 $\Phi_i$ 与 $\bar\Phi_i$ 的实解析（标量）函数，质量量纲为 2，称为凯勒势。\end{definition}

在单个手征超场 $\Phi$ 的情形：

% [原文 p.224]
\[
K(\Phi,\bar\Phi)=\sum_{m,n=1}^{\infty}c_{mn}\Phi^m\bar\Phi^n\qquad\text{其中}\quad c_{mn}=\bar c_{nm}\,,
\]
只有当对一切 $(m,n)\ne(1,1)$ 都有 $c_{mn}=0$ 时，它所定义的理论才是可重整化的，此时用 $\dd^2\theta\dd^2\bar\theta$ 积分给出 $\Phi$ 的标量与 Weyl 费米子分量的标准动能项：
\[
\int \dd^2\theta\,\dd^2\bar\theta\,\bar\Phi\Phi=\partial_\mu\bar\phi\partial^\mu\phi+\frac{\ii}{2}\left(\partial_\mu\psi\sigma^\mu\bar\psi-\psi\sigma^\mu\partial_\mu\bar\psi\right)+\bar F F+\text{全导数}\,.
\]
一般地，式 (17.10) 描述映入一个由手征超场标量分量坐标化的黎曼流形的 $\mathcal{N}=1$ 超对称 sigma 模型。SUSY 意味着该流形实际上是凯勒流形，凯勒势为 $K$（名称由此而来）。仅涉及手征超场的这类超对称 sigma 模型作为低能有效理论出现（参见 17.3 节）。靶空间是规范理论的模空间，而不可重整化性暴露了 sigma 模型描述只在低能下有效这一事实。

\begin{definition}[17.22] 对 $\Phi_i$ 的解析函数 $W$，还可以考虑 F 项：
\begin{equation}
\int \dd^2\theta\,W(\Phi_i)+\text{c.c.}=-\frac{\partial W}{\partial\phi_i}F_i-\frac{1}{2}\frac{\partial^2W}{\partial\phi_i\partial\phi_j}\psi_i\psi_j+\text{c.c.}+\text{全导数}\,,
\tag{17.11}
\end{equation}
它们导致质量项与相互作用项。$W$ 称为超势。超势项只有在 $W$ 至多为三次时才可重整化。\end{definition}

% [Tutor Note: (17.11) 的可核对特例：取 $W=\frac{m}{2}\Phi^2$，则第一项给出 $-m\varphi F$，第二项给出 $-\frac{m}{2}\psi\psi$；消去辅助场 $F$ 后应与 (17.10) 中 $|F|^2$ 项合成的标量质量项相容。又注意 Fermi 统计使 $\psi_i\psi_j$ 对 $i,j$ 对称，故二阶导数的对称化是自动的。]

总之，手征超场的 $\mathcal{N}=1$ 拉格朗日密度具有形式
\begin{equation}
\int \dd^2\theta\,\dd^2\bar\theta\,K(\Phi_i,\bar\Phi_i)+\int \dd^2\theta\,W(\Phi_i)+\int \dd^2\bar\theta\,\bar W(\bar\Phi_i)\,.
\tag{17.12}
\end{equation}

\subsubsection*{矢量超场}

\begin{definition}[17.23] 设 $V$ 是 $\mathcal{N}=1$ 矢量超场。$V$ 的超对称场强为：
\begin{equation}
W_\alpha=-\frac{1}{4}\bar D^2D_\alpha V\,,
\tag{17.13}
\end{equation}
其中右端的诸 $D$ 由式 (17.8) 定义。\end{definition}

超场场强 $W_\alpha$ 在规范变换 $V\to V+\Phi+\bar\Phi$ 下不变，并且
\begin{equation}
\int \dd^2\theta\,W^\alpha W_\alpha+\text{c.c.}=-F_{\mu\nu}F^{\mu\nu}-4\ii\lambda\sigma^\mu\partial_\mu\bar\lambda+2D^2\,,
\tag{17.14}
\end{equation}
其中 $F_{\mu\nu}=\partial_\mu v_\nu-\partial_\nu v_\mu$，并且从此以后我们省略“全导数”项。此外：

% [原文 p.225]
\[
\int \dd^2\theta\,W^\alpha W_\alpha-\text{c.c.}=\frac{\ii}{2}\epsilon^{\mu\nu\rho\sigma}F_{\mu\nu}F_{\rho\sigma}=\ii F_{\mu\nu}(\star F)^{\mu\nu}\,,
\]
其中 $\star$ 是 $4d$ 闵可夫斯基空间中的 Hodge 星算子。因此：

\begin{definition}[17.24] $4d$ 麦克斯韦理论的 $\mathcal{N}=1$ 推广，或等价地 $\mathcal{N}=1$ 矢量超场的动能（与 $\Theta$）项，为：
\begin{equation}
\frac{e^2}{4\pi}\operatorname{Im}\left(\tau_M\int \dd^2\theta\,W^\alpha W_\alpha\right)=-\frac{1}{4}F_{\mu\nu}F^{\mu\nu}-\ii\lambda\sigma^\mu\partial_\mu\bar\lambda+\frac{1}{2}D^2+\frac{\Theta e^2}{16\pi^2}F_{\mu\nu}(\star F)^{\mu\nu}\,,
\tag{17.15}
\end{equation}
其中
\begin{equation}
\tau_M=\frac{2\pi\ii}{e^2}+\frac{\Theta}{2\pi}
\tag{17.16}
\end{equation}
是复化阿贝尔规范耦合。\end{definition}

% [Tutor Note: (17.15)–(17.16) 的归一化可作核对：$\tau_M$ 的虚部 $2\pi/e^2$、实部 $\Theta/2\pi$，分别对应规范动能与 theta 项；把 $\tau_M$ 代入后实部给出 $-\frac{1}{4}F^2$（系数恰好是标准麦克斯韦归一化），虚部给出 $\frac{\Theta e^2}{16\pi^2}F(\star F)$。非阿贝尔情形 (17.19) 中把 $e$ 换成 $g$、系数变为 $4\pi i/g^2$，这与 $\tr$ 的归一化 $2\tr(T^aT^b)=\delta^{ab}$ 有关，比较时须连同 (17.18) 的 $1/16\pi$ 一并看。]

这可如下推广到非阿贝尔杨–米尔斯理论。设 $(T^a)_{a=1}^{\dim\mathfrak{g}}$ 是紧单李代数 $\mathfrak{g}$ 的厄米特生成元，其归一化使得 $2\tr(T^aT^b)=\delta^{ab}$，其中 $\tr$ 表示 $\mathfrak{g}$ 上的一个基灵型，使 $\mathfrak{g}$ 的长根的平方长度为 2。设 $V=2gV_aT^a$，其中 $V_a$ 是通常的矢量超场，$g$ 是非阿贝尔耦合常数。

\begin{definition}[17.25] 非阿贝尔 $\mathcal{N}=1$ 矢量超场的超对称场强为：
\begin{equation}
W_\alpha=-\frac{1}{4}\bar D^2\left(\ee^{-V}D_\alpha\ee^{V}\right)\,.
\tag{17.17}
\end{equation}
非阿贝尔 $\mathcal{N}=1$ 矢量超场的动能（与 $\Theta$）项写为：
\begin{equation}
\begin{aligned}
\frac{1}{16\pi}\operatorname{Im}\left(\tau\int \dd^2\theta\,\tr W^\alpha W_\alpha\right)
={}&\tr\left[-\frac{1}{2}F_{\mu\nu}F^{\mu\nu}-2\ii\lambda\sigma^\mu D_\mu\bar\lambda+D^2\right]\\
&+\frac{\Theta g^2}{16\pi^2}\tr F_{\mu\nu}(\star F)^{\mu\nu}\,,
\end{aligned}
\tag{17.18}
\end{equation}
其中
\begin{equation}
\tau=\frac{4\pi\ii}{g^2}+\frac{\Theta}{2\pi}\,,\qquad
F_{\mu\nu}=\partial_\mu v_\nu-\partial_\nu v_\mu-\ii g[v_\mu,v_\nu]\,,\qquad
D_\mu=\partial_\mu-\ii g[v_\mu,\cdot]\,,
\tag{17.19}
\end{equation}
分别是复化非阿贝尔规范耦合、非阿贝尔场强\footnote{从几何上看，反对称张量 $F_{\mu\nu}$ 最好视为取值于 $\mathfrak{g}$ 的 2-形式 $F=F_{\mu\nu}\dd x^\mu\otimes\dd x^\nu=\sum_{\mu<\nu}F_{\mu\nu}\dd x^\mu\wedge\dd x^\nu$，其中 $\dd x^\mu\wedge\dd x^\nu=(\dd x^\mu\otimes\dd x^\nu-\dd x^\nu\otimes\dd x^\mu)$。体积形式为 $\dd^4x=\dd x^0\wedge\dd x^1\wedge\dd x^2\wedge\dd x^3$。用这一记号，$\dd^4x\,F_{\mu\nu}(\star F)^{\mu\nu}=2F\wedge F$，而 theta 项可改写为
\[
k\theta g^2=\frac{\theta g^2}{16\pi^2}\int \dd^4x\,\tr F_{\mu\nu}(\star F)^{\mu\nu}=\frac{\theta g^2}{8\pi^2}\int \tr F\wedge F\,.
\]
数 $k$ 是 $F$ 为其曲率的主 $G$-丛的瞬子数。为具体起见，设 $\mathfrak{g}=\mathfrak{su}(N)$；归一化的选取使得若李代数 $\mathfrak{g}$ 对应的规范群连通且单连通，即 $G=\SU(N)$，则 $k$ 是整数 [348, 368, 181]。这意味着 $\theta$ 具有周期性 $2\pi/g^2$；对 $G=\SU(N)/\Z_k$（$\Z_k$ 是 $\SU(N)$ 的中心 $\Z_N$ 的子群）等其他规范群，角 $\theta$ 自然具有扩大的周期性，例如当 $k=N$ 时 $\theta\sim\theta+4\pi N/g^2$。瞬子将在 17.3.2 节讨论。}、以及协变导数（作用在伴随表示相联系的向量丛上）。式 (17.18) 的拉格朗日密度定义了以规范代数 $\mathfrak{g}$ 为规范代数的 $4d$ $\mathcal{N}=1$ 超杨–米尔斯（SYM）理论。\end{definition}

% [原文 p.226]
阿贝尔矢量超场 $V$ 的 D 项在（阿贝尔）规范变换下按全导数变换，从而允许拉格朗日密度中出现如下形式的附加贡献：
\begin{equation}
\xi\int \dd^2\theta\,\dd^2\bar\theta\,V(x,\theta,\bar\theta)=g\xi D(x)\,,\qquad \xi\in\R\,.
\tag{17.20}
\end{equation}
这些称为 Fayet–Iliopoulos 项。

\subsubsection*{$\mathcal{N}=1$ 规范-物质拉格朗日量}

为把规范对称性与物质耦合，需要考虑在规范群的非平凡表示 $R_i$ 中变换的手征超场 $\Phi_i$。组合变换
\begin{equation}
\ee^{V}\to\ee^{\ii\bar\Lambda_aT^a}\ee^{V}\ee^{\ii\Lambda_aT^a}\,,\qquad
\Phi_i\to\ee^{\ii\Lambda_aT^{R_i}_a}\Phi_i\,,
\tag{17.21}
\end{equation}
构成规范对称性的超对称推广，其中 $\Phi_i$ 是在规范群的表示 $R_i$ 中变换的手征超场，$T^{R_i}_a$ 是表示 $R_i$ 中的规范生成元，而 $\Lambda_a$（$a=1,\dots,\dim\mathfrak{g}$）是体现对称变换参数的手征超场。表示 $R_i$ 的 Dynkin 指标 $I(R_i)$ 定义为 $\tr_{R_i}(T^{R_i}_aT^{R_i}_b)=I(R_i)\delta^{ab}$（参见例如 [108]）。

规范不变性要求带荷手征超场的动能项修改为：
\begin{equation}
\int \dd^2\theta\,\dd^2\bar\theta\,K(\Phi_i,(\bar\Phi\ee^{2gV_R})_i)\,.
\tag{17.22}
\end{equation}
特别地，在 $G$ 的表示 $R$ 中变换的手征超场 $\Phi$ 的动能项（取标准凯勒势）成为：
\begin{equation}
\begin{aligned}
\int \dd^2\theta\,\dd^2\bar\theta\,\bar\Phi\ee^{2gV_R}\Phi={}&(D_\mu\bar\phi)D^\mu\phi-\ii\psi\sigma^\mu D_\mu\bar\psi+\bar F F\\
&+\ii\sqrt{2}g\bar\phi\lambda\psi-\ii\sqrt{2}g\bar\psi\bar\lambda\phi+g\bar\phi D\phi\,,
\end{aligned}
\tag{17.23}
\end{equation}
其中 $D_\mu=\partial_\mu-\ii gv^a_\mu T^a_R$ 是与表示 $R$ 相联系的向量丛上的协变导数（而 $D$ 是矢量多重态中的辅助场）。注意前式第二行中的每一项都含一个场

% [Tutor Note: (17.23) 第二行各项系数 $\sqrt{2}g$ 的来源：把 $e^{2gV_R}$ 展开到一阶得 $2g\,\theta\sigma^\mu\bar\theta\,v_\mu$，与 (17.9) 中 $\sqrt{2}\theta\psi$ 及 $\sqrt{2}\theta\bar\psi$ 缩并，$\dd^2\theta$ 积分给出 $\ii\sqrt{2}g$ 型耦合；核对时注意 $\theta$ 的缩并约定 $\theta\theta$ 只对 $\theta^1\theta^2$ 非零，故系数中不出现额外的 2。]

% 本页（印刷 p.226 / PDF p266）正文在此中断，续见 ch17b（p.227 起）。
